\documentclass[11pt]{article}
\usepackage[margin=1in]{geometry}
\usepackage{amsmath,amssymb}
\usepackage[hidelinks]{hyperref}

\title{Literature Status: Fixed Point Spectrum for \texorpdfstring{$L_p$}{Lp} Actions}
\author{FA/Banach run \texttt{fa\_banach\_001}}
\date{2026-06-15}

\begin{document}
\maketitle

\section*{Status}

This is a literature-already-answered packet, not an original proof from the
run.  Problems 14 and 15 in Nowak, arXiv:1302.6609, ask whether the
$L_p$ fixed point spectrum $\mathcal F(G)$ of a finitely generated
property-(T) group is connected, and whether there is a critical value
$p>2$ with
\[
  \mathcal F(G)=(1,p).
\]
Marrakchi--de la Salle, arXiv:2001.02490, answer these questions by proving
that the complementary set of exponents admitting affine isometric
$L_p$ actions with unbounded orbits is an upper interval.

\section*{Original Problem}

The original source is:

\begin{quote}
P. W. Nowak, \emph{Group Actions on Banach Spaces}, arXiv:1302.6609.
\end{quote}

On page 10, Nowak defines
\[
\mathcal F(G)=\{p\in(1,\infty): H^1(G,\pi)=0
\text{ for every isometric representation }\pi\text{ on an }L_p\text{-space}\}.
\]
The paper then asks:

\begin{itemize}
\item Problem 14 (C. Drutu): is $\mathcal F(G)$ connected?
\item Problem 15: given $G$ with property (T), is there a critical value
      $p>2$ such that $\mathcal F(G)=(1,p)$?
\end{itemize}

\section*{Later Answer}

The supporting source is:

\begin{quote}
A. Marrakchi and M. de la Salle,
\emph{Isometric actions on $L_p$-spaces: dependence on the value of $p$},
arXiv:2001.02490.
\end{quote}

Their Theorem 1 shows monotonicity in $p$ for affine isometric actions on
$L_p$ spaces.  Corollary 2 then states that, for every topological group $G$,
the set of $p\in(0,\infty)$ for which $G$ admits a continuous affine
isometric action on an $L_p$ space with unbounded orbits has the form
\[
  (p_G,\infty) \quad\text{or}\quad [p_G,\infty)
\]
for some $p_G\in\{0\}\cup[2,\infty]$.  The same page says this answers the
question raised by Chatterji--Drutu--Haglund, often referred to as Drutu's
conjecture.

For second countable locally compact groups with property (T), the supporting
paper records the sharper endpoint information: $p_G>2$, and there is an
affine isometric $L_{p_G}$ action without a fixed point.  Finitely generated
groups are second countable locally compact in the discrete topology, so this
applies to Nowak's stated setting.

\section*{Translation to Nowak's Spectrum}

For $1<p<\infty$, $L_p$ spaces are reflexive.  In this range, an affine
isometric action has a fixed point if and only if its orbits are bounded; this
is recalled on page 2 of arXiv:2001.02490 and is also the fixed-point
principle used in Nowak's source.  Therefore
\[
  p\in\mathcal F(G)
  \quad\Longleftrightarrow\quad
  \text{there is no affine isometric }L_p\text{ action with unbounded orbits}.
\]

Since the unbounded-orbit set is $[p_G,\infty)$ for property-(T) groups in
Nowak's class, its complement in $(1,\infty)$ is exactly
\[
  \mathcal F(G)=(1,p_G),
\]
with $p_G>2$.  This proves connectedness and the critical-value form asked
for in Problems 14 and 15.

\section*{Verification Notes}

The answer is not contained in arXiv:1302.6609; that paper explicitly says
the problems were far from understood.  The later paper is independent
literature and explicitly identifies Corollary 2 as an answer to Drutu's
conjecture.  The only mathematical translation needed is the standard
equivalence, for $1<p<\infty$, between vanishing first cohomology for all
isometric representations and fixed points for all affine isometric actions
on $L_p$ spaces.

\section*{Limitations}

This packet does not compute the numerical value of $p_G$ for any specific
group.  It also does not address the other questions in Nowak's survey,
including the $\widetilde A_2$ question, uniformly bounded Hilbert
representations, reduced $\ell_p$ cohomology, Yu's question, or Nica's
mapping-class-group conjecture.

\section*{Human Review Recommendation}

Verify the endpoint statement on page 2 of arXiv:2001.02490 and the
fixed-point/bounded-orbit translation for affine isometric actions on
$L_p$, $1<p<\infty$.  These two points are the complete bridge from the later
paper to Nowak's Problems 14 and 15.

\end{document}
